\chapter{Open Gaps}
\label{sec:origins-gaps}
% ===================================================================

\begin{enumerate}[label=\textbf{Gap \arabic*.}, leftmargin=*, itemsep=8pt]

  \item \textbf{Density of replicating terms
        (Conjecture~\ref{conj:density}).}
        The argument that the added replication machinery can always
        be hidden from any given HML formula requires a careful
        analysis of which formulae can probe the replication channel.
        In particular, if $\phi = \langle K \rangle \psi$ where $K$
        involves the replication channel, the argument fails.
        A precise density theorem would need to show that such
        formulae form a meager (first-category) set in the logical
        topology, so that ``generically'' $\phi$ does not probe
        replication.

  \item \textbf{Defining the elementary terms.}
        Proposition~\ref{prop:assembly-index} requires a designated
        set $E$ of elementary terms.
        In the chemical setting, these are atoms or monomers.
        In the Rho calculus, the natural candidates are the
        behavioral primes of Part~II (Section~4): the irreducible
        terms from which all others are built by interaction.
        Confirming that this identification gives assembly indices
        matching those of Assembly Theory requires working through
        the combinatorics of the Rho calculus prime factorization.

  \item \textbf{The basin of attraction.}
        The claim that a population entering a neighborhood of
        $\mathrm{Rep}$ is ``sucked in'' to the replicating fixed
        point requires a Lyapunov-style argument: a function on the
        process space that is monotonically decreasing along
        trajectories near the fixed point.
        The natural candidate is the distance $d_{\HML}$ to
        $\mathrm{Rep}$, but showing it decreases requires
        assumptions on the interaction dynamics.

  \item \textbf{The compact cluster as the new geosphere.}
        The identification of Smith's fourth geosphere with the
        compact coherence cluster centered on the replicating fixed
        point requires showing that this cluster is indeed compact
        (in the sense of Part~II) and that it persists under the
        perturbations of the selection dynamics.

  \item \textbf{Crossover as a variant of comm.}
        The claim that genetic recombination is one definitional
        step from the comm rule requires specifying precisely what
        ``fragment substitution'' means in the Rho calculus and
        showing that the resulting operator has the recombination
        properties of a genetic crossover (fitness inheritance,
        schema theorem, etc.).

  \item \textbf{Gradient-coupled rules and the weight map.}
        Definition of gradient-coupled rules modifies the weight
        map of Part~I by introducing an external parameter $\Phi$.
        A precise treatment requires extending the weighted GSLT
        framework to allow time-varying or externally-driven weight
        maps, and showing that the path integral remains
        well-defined under this extension.

  \item \textbf{Metabolic viability (Conjecture~\ref{conj:viable}).}
        The condition that catabolic inflow exceeds anabolic cost
        in expectation requires computing the path integral of a
        driven open system --- a substantially harder problem than
        the closed-system path integral of Part~I.
        The quantum Gillespie algorithm of Part~I provides a
        simulation method, but an analytic viability criterion
        requires further development.

  \item \textbf{Trophic level as causal depth
        (Proposition~\ref{prop:trophic}).}
        The identification of trophic level with
        $d^o_{\min}(K_\Phi, B)$ requires showing that the
        account-flow graph $\mathcal{F}$ is faithfully captured
        by the open synchronization tree structure.
        In particular, it requires that every account transfer
        corresponds to a minimal-context transition, which is
        plausible but needs verification.

\end{enumerate}

% ===================================================================
